\documentclass[12pt,a4paper]{article}
\usepackage[utf8]{inputenc}
\usepackage[spanish]{babel}
\usepackage{geometry}
\geometry{top=2.5cm, bottom=2.5cm, left=3cm, right=3cm}
\usepackage{graphicx}
\usepackage{amsmath}
\usepackage{booktabs}
\usepackage{hyperref}
\usepackage{float}
\usepackage{caption}
\usepackage{fancyhdr}
\usepackage{xcolor}

% Configuration of hyperref
\hypersetup{
    colorlinks=true,
    linkcolor=cyan,
    filecolor=magenta,      
    urlcolor=blue,
}

% Page headers and footers
\pagestyle{fancy}
\fancyhf{}
\rhead{\textcolor{gray}{Proyecto Académico Sirocco V1}}
\lhead{\textcolor{gray}{Diseño de Turbomáquinas}}
\cfoot{\thepage}
\newcommand{\hrules}{\rule{\linewidth}{0.5mm}}

\begin{document}

% --- PORTADA ---
\begin{titlepage}
    \centering
    \vspace*{1cm}
    {\scshape\LARGE Universidad Nacional\par}
    \vspace{1.5cm}
    {\scshape\Large Facultad de Ingeniería\par}
    \vspace{0.5cm}
    {\large Curso: Diseño de Turbomáquinas\par}
    \vspace{2cm}
    \hrules
    \vspace{0.4cm}
    {\huge\bfseries Diseño y Cálculo Reproducible de un Ventilador Centrífugo Tipo Sirocco V1\par}
    \vspace{0.4cm}
    \hrules
    \vspace{1.5cm}
    {\large\bfseries Grupo 2\par}
    \vspace{0.5cm}
    {\large Revisión: Versión V1 (Cierre Unidimensional)\par}
    \vfill
    {\large \today\par}
\end{titlepage}

\tableofcontents
\newpage

% --- SECCIÓN 1: INTRODUCCIÓN ---
\section{Introducción y Alcance}
El presente documento expone la memoria técnica y el predimensionamiento de un ventilador centrífugo multialabe del tipo Sirocco, de doble ancho y doble entrada (DIDW), diseñado para operar con aire limpio en condiciones normales.

El punto de diseño especificado para el proyecto es el siguiente:
\begin{itemize}
    \item \textbf{Caudal de diseño ($Q$):} $3.0\text{ m}^3/\text{s}$ ($10,800\text{ m}^3/\text{h}$).
    \item \textbf{Presión estática del ventilador ($\Delta p_s$):} $90\text{ mmH}_2\text{O}$ ($882.6\text{ Pa}$).
    \item \textbf{Densidad de diseño del aire ($\rho$):} $1.20\text{ kg/m}^3$.
\end{itemize}

El dimensionamiento actual constituye la fase de análisis unidimensional (Revisión V1). Los resultados que se muestran a continuación establecen la geometría base adecuada para el posterior modelado CAD y la validación aerodinámica tridimensional mediante dinámica de fluidos computacional (CFD).

% --- SECCIÓN 2: CONFIGURACIÓN ---
\section{Configuración Seleccionada}
Debido al elevado caudal requerido ($3.0\text{ m}^3/\text{s}$), un ventilador de entrada simple (SISW) exigiría un diámetro de succión excesivamente grande, lo que derivaría en altas pérdidas secundarias en la boca de entrada o velocidades axiales muy elevadas.

Para mitigar este problema, se ha seleccionado una configuración de \textbf{Doble Entrada y Doble Ancho (DIDW)}. Esta arquitectura divide el caudal de succión en dos ojos de entrada iguales ($1.5\text{ m}^3/\text{s}$ cada uno), reduciendo significativamente la velocidad axial en el ojo a un valor óptimo de aproximadamente $7.77\text{ m/s}$, atenuando el ruido tonal y las pérdidas aerodinámicas por choque en la entrada de la corona.

% --- SECCIÓN 3: GEOMETRÍA ---
\section{Parámetros Geométricos del Rodete}
La geometría calculada para la revisión V1 del rodete Sirocco presenta las siguientes dimensiones y coeficientes constructivos:

\begin{table}[H]
\centering
\caption{Dimensiones del Rodete Sirocco V1}
\begin{tabular}{lcr}
\toprule
\textbf{Parámetro} & \textbf{Símbolo} & \textbf{Valor} \\
\midrule
Diámetro exterior del rodete & $D_2$ & $600\text{ mm}$ \\
Diámetro de entrada del rodete & $D_1$ & $510\text{ mm}$ \\
Relación de diámetros & $D_1/D_2$ & $0.850$ \\
Ancho constructivo del rodete & $b$ & $230\text{ mm}$ \\
Número de álabes & $Z$ & $48$ \\
Ángulo de salida del álabe & $\beta_2$ & $125^\circ$ \\
Ángulo de entrada del álabe & $\beta_1$ & $18^\circ$ \\
Espesor del álabe & $t$ & $1.2\text{ mm}$ \\
\bottomrule
\end{tabular}
\end{table}

La relación de diámetros de $0.85$ es la estándar para ventiladores Sirocco. Permite álabes de longitud radial corta, dejando un espacio central amplio para distribuir el aire axial. El ángulo $\beta_2 = 125^\circ$ indica álabes fuertemente curvados hacia adelante, diseñados para transferir una gran componente cinética al flujo.

% --- SECCIÓN 4: CÁLCULOS ---
\section{Cálculos Fluidodinámicos y Velocidades}
A continuación, se detallan los resultados del análisis del triángulo de velocidades a la salida del rodete a la velocidad recomendada de operación de \textbf{1100 RPM}:

\subsection{Velocidades de la Corriente a la Salida}
\begin{itemize}
    \item \textbf{Velocidad Periférica ($U_2$):} 
    \begin{equation}
        U_2 = \frac{\pi \cdot D_2 \cdot n}{60} = 34.56\text{ m/s}
    \end{equation}
    \item \textbf{Velocidad Radial/Meridional ($C_{m2}$):} Calculada descontando el bloqueo efectivo por espesor de los álabes ($k_{b2} = 0.963$):
    \begin{equation}
        C_{m2} = \frac{Q}{\pi \cdot D_2 \cdot b \cdot k_{b2}} = 7.19\text{ m/s}
    \end{equation}
    \item \textbf{Factor de Deslizamiento (Wiesner):} Permite corregir el ángulo real de salida del aire debido al guiado finito:
    \begin{equation}
        \sigma = 1 - \frac{\sqrt{\sin\beta_2}}{Z^{0.7}} = 0.940
    \end{equation}
    \item \textbf{Componente Tangencial ($C_{u2}$):}
    \begin{equation}
        C_{u2} = \sigma \cdot U_2 - C_{m2} \cdot \cot\beta_2 = 37.51\text{ m/s}
    \end{equation}
\end{itemize}

\subsection{Potencia y Eficiencia}
Con estas velocidades, el incremento de presión teórico de Euler es de $\Delta p_E = \rho \cdot U_2 \cdot C_{u2} = 1555.5\text{ Pa}$. 

La potencia al eje requerida (asumiendo una eficiencia de transmisión mecánica del $96\%$) se calcula en \textbf{4.86 kW}, lo que arroja una eficiencia estática estimada del \textbf{54.5\%} para el punto de operación. Por lo tanto, se ha prescrito la selección de un motor de \textbf{5.5 kW} (con un margen de seguridad del $13.1\%$).

% --- SECCIÓN 5: VOLUTA ---
\section{Geometría de la Voluta}
El aire abandona la salida del rodete con una velocidad absoluta muy elevada ($C_2 \approx 38.2\text{ m/s}$). La carcasa espiral (voluta) tiene la función de actuar como difusor para reconvertir esta energía cinética en presión estática:

\begin{itemize}
    \item \textbf{Ancho interior de la voluta:} $280\text{ mm}$, holgura lateral adecuada respecto al rodete de $230\text{ mm}$ para prevenir fricciones por desalineación.
    \item \textbf{Holgura de la Lengua (Cut-off):} Se sitúa a $48\text{ mm}$ del rodete ($8\%$ de $D_2$). Esta holgura controla el nivel de ruido aerodinámico provocado por las pulsaciones periódicas de paso de los álabes.
    \item \textbf{Velocidad de descarga final:} $\approx 15\text{ m/s}$ en la boca de salida de la voluta.
\end{itemize}

% --- SECCIÓN 6: COMPARATIVA ---
\section{Análisis Operativo: 1100 RPM vs 1500 RPM}
Durante los tanteos preliminares se evaluó una velocidad de rotación de $1500\text{ RPM}$. Sin embargo, se descartó debido a las siguientes restricciones:

\begin{table}[H]
\centering
\caption{Comparativa de Velocidades de Giro}
\begin{tabular}{lccl}
\toprule
\textbf{Magnitud} & \textbf{1100 RPM (V1)} & \textbf{1500 RPM (Tanteo)} & \textbf{Análisis} \\
\midrule
Velocidad Periférica ($U_2$) & $34.56\text{ m/s}$ & $47.12\text{ m/s}$ & Mayor esfuerzo estructural \\
Presión de Euler ($\Delta p_E$) & $1555.5\text{ Pa}$ & $2892.4\text{ Pa}$ & Exceso innecesario de presión \\
Potencia al Eje Estimada & $4.86\text{ kW}$ & $9.03\text{ kW}$ & Excede motor de $5.5\text{ kW}$ \\
Estado de Viabilidad & \textbf{VIABLE} & \textbf{INVIABLE} & Seleccionada a 1100 RPM \\
\bottomrule
\end{tabular}
\end{table}

A 1500 RPM la potencia demandada ($9.03\text{ kW}$) superaba con creces el límite del motor de $5.5\text{ kW}$ disponible comercialmente para esta gama de tamaño del ventilador. La selección de 1100 RPM ajusta la presión útil al punto requerido con un consumo de potencia óptimo.

% --- SECCIÓN 7: REPOSITORIO ---
\section{Estructura del Proyecto Académico}
El repositorio oficial del proyecto (\url{https://github.com/Aquiles-ryp360/diseno-ventilador-sirocco}) se encuentra organizado bajo las siguientes carpetas:
\begin{itemize}
    \item \textbf{00\_gestion:} Planificación y cronogramas.
    \item \textbf{04\_bibliografia:} Literatura y catálogos técnicos consultados.
    \item \textbf{05\_avances:} Informes periódicos entregados al docente.
    \item \textbf{06\_calculos:} Memorias de cálculo en formato Markdown.
    \item \textbf{07\_planos:} Modelos CAD y planos DXF de fabricación.
    \item \textbf{08\_software:} Scripts de cálculo desarrollados en Python y Octave.
    \item \textbf{presentacion:} Diapositivas web interactivas (HTML/CSS/JS) y formato PowerPoint.
\end{itemize}

\end{document}
